\subsection{Overview}

Figure~\ref{fig:model} summarizes the venue congruence model. Its starting point is an employee's perception that the employer has broken an ethical commitment, which we call ideological contract breach. The model proposes that such a breach is more likely to be felt as a violation when the commitment was learned through the employer's public communication, and that the same feature of the promise, its publicity, pushes the employee's response toward public rather than internal venues. Internal voice safety sets the boundary: where raising concerns inside feels safe and worthwhile, the internal channel is expected to absorb much of the response; where it does not, the response is expected to go public or, when the promise was private, to turn into exit or silence. The final proposition concerns Generation Z and is stated as two rival accounts. Every arrow in the figure is a proposition, and none has been tested.

\begin{figure}[htb]\centering
\resizebox{\linewidth}{!}{%
\begin{tikzpicture}[font=\small, box/.style={draw, rounded corners=2pt, align=center, minimum height=9mm, text width=27mm, fill=white}, mod/.style={draw, dashed, rounded corners=2pt, align=center, minimum height=8mm, text width=24mm, fill=white}, resp/.style={draw, align=center, minimum height=7mm, text width=24mm, fill=white}, grp/.style={draw, gray, rounded corners=3pt, inner sep=2.5mm}, arr/.style={-{Stealth[length=2.2mm]}, thick}, marr/.style={-{Stealth[length=2mm]}, dashed}, rarr/.style={-{Stealth[length=2mm]}, dotted}]
\node[box] (icb) at (0,0) {Ideological\\contract breach};
\node[box] (fev) at (4.6,0) {Felt ethical\\violation};
\node[resp] (ivo) at (12.6,3.0) {Internal ethical voice};
\node[resp] (pvo) at (12.6,1.75) {Public ethical voice};
\node[resp] (ext) at (12.6,-1.75) {Exit};
\node[resp] (sil) at (12.6,-3.0) {Silence};
\node[grp, fit=(ivo)(pvo), label={[font=\footnotesize, gray]above:Venue of ethical voice}] (venue) {};
\node[grp, fit=(ext)(sil), label={[font=\footnotesize, gray]below:Withdrawal}] (wd) {};
\node[mod] (pp) at (2.3,3.4) {Promise\\publicity};
\node[mod] (ivs) at (8.9,0) {Internal voice\\safety};
\node[draw, dotted, rounded corners=2pt, align=center, minimum height=8mm, fill=white] (rival) at (1.6,-3.2) {Rival accounts of the\\Generation Z pattern\\P5a: cohort\\P5b: career stage};
\draw[arr] (icb.east) -- (fev.west);
\node[font=\footnotesize, below] at (2.3,0) {P1};
\coordinate (vstart) at ($(fev.east)+(0,0.2)$);
\coordinate (wstart) at ($(fev.east)+(0,-0.2)$);
\draw[arr] (vstart) -- (venue.west);
\draw[arr] (wstart) -- (wd.west);
\draw[marr] (pp.south) -- (2.3,0.05);
\path (vstart) -- (venue.west) coordinate[pos=0.74] (p2pt);
\draw[marr] (pp.east) -| node[pos=0.25, above, font=\footnotesize] {P2} (p2pt);
\path (vstart) -- (venue.west) coordinate[pos=0.55] (p3pt);
\path (wstart) -- (wd.west) coordinate[pos=0.55] (p4pt);
\draw[marr] (ivs.north) -- node[right, font=\footnotesize] {P3} (ivs.north |- p3pt);
\draw[marr] (ivs.south) -- node[right, font=\footnotesize] {P4} (ivs.south |- p4pt);
\draw[rarr] (rival.west) .. controls +(-2.6,0) and +(-2.6,0) .. (pp.west);
\draw[rarr] (rival.east) .. controls +(4.2,0) and +(-0.6,-1.2) .. (ivs.south west);
\end{tikzpicture}}
\caption{The venue congruence model. Solid arrows show the proposed paths from breach to felt violation and from felt violation to the venue of ethical voice or to withdrawal. Dashed arrows show proposed moderation: P1 (promise publicity on the breach to violation path), P2 (promise publicity shifting the venue toward public voice), P3 (internal voice safety shifting the venue toward internal voice), and P4 (low internal voice safety, combined with low promise publicity, directing violation toward exit and silence). Dotted arrows show the constructs through which the rival accounts in Proposition 5 operate. All paths are propositions awaiting empirical test.}\label{fig:model}
\end{figure}

\subsection{Constructs}

Ideological contract breach is the employee's perception that the organization has failed to honor a commitment to an ethical principle or valued cause that the employee believed it had made. The definition follows~\citet{thompson2003violations}, narrowed to commitments with ethical content, such as fair treatment of workers in the supply chain, honest dealing with customers, or environmental conduct. It is a perception, and as~\citet{morrison1997employees} stress, a perceived breach can arise from a real failure to deliver or from a mismatch between what each party understood the promise to be.

Promise publicity is our own construct, and no existing measure captures it. We define it as the extent to which an employee's understanding of the employer's ethical obligations was formed through communication addressed to the public, such as employer branding, corporate responsibility statements, and the organization's social media, rather than through communication addressed to the employee alone, such as a recruiter's assurance or a manager's private commitment. The construct builds on~\citet{rousseau2001schema}, who located contract formation partly in pre-entry schemas and in promises conveyed through many channels, and on employer branding research, which describes firms deliberately promoting an employer image both inside and outside the organization~\citep{backhaus2004conceptualizing}. It differs from perceived corporate social performance, which concerns how good the employer's conduct is; promise publicity concerns where the employee learned what the employer had promised.

Felt ethical violation is the emotional response of anger and betrayal that can follow a perceived breach~\citep{morrison1997employees}, applied here to ethical commitments. Meta-analytic evidence links breach to violation and treats violation as one of the affective states through which breach affects attitudes and behavior~\citep{zhao2007impact}.

Internal voice safety is the employee's belief that raising an ethical concern through channels inside the organization, such as a supervisor, a senior manager, or an ethics function, will be safe and has a reasonable chance of being acted on. It adapts psychological safety, a shared belief that a team is safe for interpersonal risk taking~\citep{edmondson1999psychological}, to the specific act of ethical voice, and it adds the efficacy judgment that voice research treats as a separate consideration from safety~\citep{morrison2014employee}.

The four responses are internal ethical voice (raising the concern with someone inside the organization who could act on it), public ethical voice (disclosing or criticizing the conduct to an audience outside the organization, including through social media or the press), exit (leaving or actively preparing to leave), and silence (withholding the concern while staying). The internal and public categories correspond to the internal and external channels distinguished in whistleblowing research~\citep{dworkin1998internal,kaptein2011inaction}. Silence is treated in the sense developed by~\citet{vandyne2003conceptualizing}, as a motivated withholding rather than an absence of opinion.

A note on design choices. We define the response as four categories rather than as a single voice variable because the model's central claim concerns substitution between venues; a single voice score would hide that. We define internal voice safety at the level of ethical voice rather than using general psychological safety because an employee can feel safe proposing a process improvement and still feel exposed raising an ethical objection, a distinction consistent with evidence that safety relates more strongly to prohibitive than to promotive voice~\citep{liang2012psychological}. Finally, we do not model loyalty as a separate response. In the organizational version of the exit-voice-loyalty framework, loyalty is a passive but constructive response of waiting for conditions to improve~\citep{farrell1983exit}; an employee who stays quiet in that hope falls within our silence category, and the motive-based distinctions of~\citet{vandyne2003conceptualizing} allow such cases to be separated in measurement.

\subsection{Breach, Promise Publicity, and Felt Ethical Violation}

In the account of~\citet{morrison1997employees}, a perceived breach can stem from reneging, where the organization knowingly fails to deliver, or from incongruence, where the two sides understood the promise differently, and whether the breach turns into violation depends on how the employee interprets it, including judgments about why it happened. \citet{robinson2000development} found in a longitudinal study of managers that breach was more strongly associated with violation when employees attributed the breach to purposeful reneging and felt unfairly treated in how it was handled.

Our argument, which is the paper's own and has not been tested, is that promise publicity makes the reneging interpretation more available. A commitment stated publicly is usually repeated and can be checked. The employee can return to the website, the report, or the campaign and confirm that the organization did say what it said. That leaves less room for the thought that the employee misunderstood, which is the thought that supports an incongruence reading. A privately communicated commitment, by contrast, is easier to reinterpret as a misunderstanding or as the view of one recruiter rather than the organization.

Public statements are not always specific, however. Aspirational wording about caring for people or the planet gives an employee little to check, and a breach of such wording may be easier to dismiss as a misunderstanding. The argument therefore applies with most force to public commitments specific enough to be verified, and specificity is a condition that tests of the first proposition should measure.

Public commitments also invite the charge of hypocrisy in a way private ones do not. Institutional research has documented how organizations adopt ethics structures that are easily decoupled from everyday practice, particularly when the pressure to adopt them comes from outside~\citep{weaver1999integrated}. Later work has argued that decoupling remains common, though increasingly as a gap between the practices organizations adopt and the ends those practices are supposed to serve~\citep{bromley2012smoke}. Either kind of gap can surface as a broken ethical promise. An employee who discovers one after joining on the strength of the public claim has reason to read the breach as deliberate rather than accidental.

\begin{proposition}
Ideological contract breach is positively related to felt ethical violation, and this relationship is stronger when promise publicity is high than when it is low.
\end{proposition}

The main effect in this proposition restates, for ethical commitments, an association already supported for psychological contracts in general~\citep{zhao2007impact}. The moderation by promise publicity is new and untested. In the propositions that follow we treat felt violation as the proximal driver of the employee's response, in line with meta-analytic models that place violation between breach and its outcomes~\citep{zhao2007impact}, and we do not propose separate direct paths from breach to the responses.

\subsection{Venue Congruence}

The second proposition is the core of the model. Once violation is felt, the employee must decide what to do with it. We propose that the venue in which the promise was made shapes the venue in which its breach is contested. We call this venue congruence. The reasoning is ours, drawing on two established ideas.

The first is that a public promise has a public audience. When an organization states an ethical commitment to customers, investors, and prospective employees at large, the employee who later sees it broken is one member of that audience among many. Raising the matter publicly then has a logic that it lacks when the promise was made privately: the employee is holding the organization to account before the audience to whom the promise was made. Research on social activism distinguishes insiders from outsiders who press organizations to change and examines how claims and pressure move between the two~\citep{briscoe2016social}. An employee who sees a public promise broken may think of disclosure less as disloyalty and more as reminding the audience of what it was told. We know of no study that has tested whether employees reason this way.

The second idea concerns the channel itself. Social media make an employee's communication visible to large audiences and keep it available over time~\citep{treem2013social}, and practitioner-oriented work has argued that these tools give employees a voice that can reach far beyond the organization, with potential to help or harm its reputation~\citep{miles2014employee}. When the original promise lives in the same public space, answering it there requires no change of setting. Evidence from whistleblowing research shows that external and internal whistleblowers differ in their circumstances and in the retaliation they face~\citep{dworkin1998internal}, which supports treating venue as an outcome in its own right, although that research does not address the origin of the promise.

\begin{proposition}
Among employees who experience felt ethical violation, promise publicity is positively related to the likelihood of public ethical voice relative to internal ethical voice.
\end{proposition}

This proposition predicts a relative effect. It does not claim that high promise publicity makes internal voice rare; it claims that the balance between the two venues shifts.

\subsection{Internal Voice Safety as a Boundary Condition}

Venue congruence describes a pull toward public voice. Whether that pull prevails should depend on what the internal channel offers. Voice research has shown repeatedly that employees weigh the personal risks and the likely usefulness of speaking up before doing so~\citep{morrison2014employee,morrison2023employee}. Fear of being labeled negatively and of damaging relationships is a leading reason employees give for staying silent~\citep{milliken2003exploratory}, and fear more broadly has been theorized as a central driver of workplace silence~\citep{kishgephart2009silenced}. Employees also hold implicit theories about when speaking up is risky, and these theories predict silence beyond what other factors explain~\citep{detert2011implicit}. On the other side, managerial openness is associated with voice, and the association runs through employees' sense of psychological safety~\citep{detert2007leadership}.

Evidence specific to ethical concerns points the same way. In a study of employee responses to observed wrongdoing, several dimensions of an organization's ethical culture were negatively related to intended external whistleblowing and inaction and positively related to intended internal responses such as reporting to management~\citep{kaptein2011inaction}. That finding concerns ethical culture rather than voice safety as we define it, and it did not examine promise publicity, but it is consistent with the idea that a credible internal channel draws responses inward.

\begin{proposition}
Internal voice safety moderates the relationship between felt ethical violation and the venue of ethical voice: when internal voice safety is high, felt ethical violation is more strongly related to internal ethical voice and less strongly related to public ethical voice than when internal voice safety is low.
\end{proposition}

Taken together with Proposition 2, this implies that public ethical voice should be most likely when promise publicity is high and internal voice safety is low. We state the two-way propositions separately to keep each testable on its own.

\subsection{When Neither Venue Seems Viable}

Some employees who feel violated will have neither a safe internal channel nor a reason, rooted in the promise, to go public. For them, the model returns to Hirschman's starting point: when voice appears futile, exit becomes the main alternative~\citep{hirschman1970exit}. The exit-voice-loyalty literature adds that the pull of exit depends on the quality of alternatives~\citep{rusbult1988impact}, and~\citet{turnley1999impact} reported that the availability of attractive alternatives moderated the link between contract violation and exit. Where exit is costly, the remaining option is to stay and say nothing, which voice research describes as defensive or acquiescent silence depending on the motive~\citep{vandyne2003conceptualizing}.

\begin{proposition}
When internal voice safety is low and promise publicity is low, felt ethical violation is more strongly related to exit and to silence than to public ethical voice.
\end{proposition}

This proposition matters for how organizations read the absence of complaints. If it holds, a quiet workforce after an ethical breach is not evidence that the breach went unnoticed; some of the people who noticed may simply have left or stopped speaking.

\subsection{Generation Z: Cohort or Career Stage?}

Why bring Generation Z into this model at all? Because the pattern the model predicts, a willingness to go public combined with reluctance to speak internally, is one that practitioner commentary has attributed to the youngest workers. A practitioner survey in the United Kingdom reported that its youngest respondents were less likely than older ones to say they would raise concerns with their employer and more likely to say they would post about workplace problems on social media or approach journalists~\citep{protect2025talkin}. That survey was not peer reviewed, measured stated intentions rather than behavior, compared age groups at a single point in time, and rested on a small youngest-age subsample. It is a reason to ask the question, not evidence for an answer.

Two explanations fit such a pattern, and the model can express both. The first is a cohort account. People who entered the labor market after public employer communication and social media became pervasive may form more of their ethical expectations of employers from public sources, raising promise publicity, and may find public channels more familiar. Both claims are our conjectures; we are not aware of peer-reviewed evidence that establishes either for Generation Z specifically, and reviews caution that cohort differences in work attitudes and values have generally been small or inconsistent~\citep{costanza2012generational,parry2011generational,rudolph2021generations}. If the cohort account is right, cohort membership should matter even among employees at the same career stage.

The second is a career-stage account. Younger workers are, almost by definition, more likely than older ones to be newcomers in junior roles. Newcomers are still establishing role clarity and social acceptance~\citep{bauer2007newcomer}. Fear of the personal costs of speaking up is a common reason for silence~\citep{milliken2003exploratory,kishgephart2009silenced}, and it is plausible, though we know of no direct test for ethical concerns, that these costs weigh more heavily on employees with little tenure or status. In a study of wrongfully dismissed whistleblowers, those who went outside the organization had shorter tenure than those who reported internally~\citep{dworkin1998internal}. On this account, low internal voice safety, not cohort, produces the pattern, and older newcomers should behave much like younger ones.

\begin{proposition}
(a) Cohort account: holding organizational tenure and hierarchical position constant, members of Generation Z report higher promise publicity and show a stronger venue congruence effect than members of older cohorts. (b) Career-stage account: once organizational tenure and hierarchical position are held constant, differences between Generation Z and older cohorts in public ethical voice disappear, and the apparent generational pattern is accounted for by the lower internal voice safety associated with newcomer and junior status.
\end{proposition}

The two parts make opposing predictions, so evidence that supports one counts against the other. One caution applies to the cohort account. At any single point in time cohort and age move together, so a study conducted once can separate membership in the youngest group from tenure and position, but it cannot tell whether a membership effect reflects cohort or age. That question needs the repeated designs discussed in the next section. We do not favor either in advance. This framing takes seriously the argument of critics that differences attributed to generations are often better explained by age, period, or developmental and career processes~\citep{rudolph2017considering,rudolph2021generations}, while leaving room for a cohort effect if the evidence supports one.
